During my PhD, Thomas Lehéricy and I studied normal approximation for sums
of random variables with a dependency graph.
In our paper~\cite[Lemma~5.1]{JL}, we sharpened a cumulant
estimate that my advisor, Ashkan Nikeghbali, had developed with Valentin
Féray and Pierre-Loïc Méliot in their work on mod-$\phi$
convergence~\cite[Theorem~9.8]{FMN}. This gave us bounds on the Kolmogorov
distance between the standardized sum and the standard normal distribution.
For some ranges of the maximum degree $D$, our bounds improved on the
Berry--Esseen bound of Chen and Shao~\cite[Theorem~2.7]{ChenShao}.
For bounded summands our Fourier argument gave an essentially optimal
estimate. With only a moment of order $2<\delta\le 3$, it gave a slower
rate~\cite[Theorems~2.1 and~2.4 and Section~4]{JL}.

A simple example motivated our conjecture. Suppose the graph consists of
disjoint cliques of size $D+1$, with identical summands in each clique.
Distinct cliques are independent by the definition of a dependency graph.
After grouping, $N$ summands give just $N/(D+1)$ independent terms. We
conjectured that the Berry--Esseen estimate for independent random variables
would still hold for every dependency graph of maximum degree $D$, with
$N/(D+1)$ in place of $N$ in the rate~\cite[Conjecture~1.4]{JL}.

Martin Raič had already conjectured the case $\delta=3$ in
2003~\cite[Conjecture~5.4]{Raic}. Just before stating his conjecture, he
suggested that the induction used in the independent case might also work
for local dependence. For $m$-dependent sequences, Shergin had proved the
corresponding moment bound in 1979~\cite[Theorem~2]{Shergin}. For general
dependency graphs, Cai, Shao and Zhang~\cite[Theorem~3.1]{CSZ} recently
obtained a bound under a fourth-moment assumption.

I finished my PhD without a proof of the conjecture. As AI models improved,
I kept giving them the problem. I first obtained a complete proof with
OpenAI's GPT-6.2~Astra in November 2026. The complete Lean~4 formalization
and instructions for checking the proof are included in the arXiv ancillary files.

The proof gives a stronger result. Let $W=\sum_i X_i$ have variance one,
where the $X_i$ are centered and square-integrable and have a dependency
graph of maximum degree $D$. Write $q=D+1$, let $Z$ be standard normal,
and denote Kolmogorov distance by $\dK$. Then
\[
 \dK(W,Z)\le 250{,}000\,q\sum_i
 \E\!\left[X_i^2\min\{1,q|X_i|\}\right].
\]
At $q=1$ this recovers the truncated-moment form of Osipov's inequality
for independent summands~\cite{Osipov}. Since $t^2\min\{1,t\}\le t^\delta$
for every $t\ge0$ and $2<\delta\le3$, the same constant gives
$\dK(W,Z)\le 250{,}000\,q^{\delta-1}\sum_i\E|X_i|^\delta$ throughout
the conjectured range. Theorem~\ref{thm:published} states the result in
the notation of our original conjecture.

\paragraph{Idea of the proof.}
We use Stein's method~\cite{Stein72} and induct on the number of summands,
as in the independent case~\cite{Bolthausen,Raic}. Smoothing the indicator
of a half-line in Stein's equation gives a sum over adjacent pairs
$(i,j)$, including $i=j$. Each term asks whether $W$ lies in a short
interval. The difficulty is that its width depends on all the summands
in the neighborhoods of $i$ and $j$. Chen and Shao estimated interval
probabilities under local dependence by concentration inequalities~\cite[Section~3]{ChenShao}.

Lemma~\ref{lem:charge} is the step that makes induction possible.
Its proof uses the fact that neighbors carrying at least half the average
weight in a closed neighborhood account for at least half its total weight.
After summing over pairs, the lemma gives intervals whose
widths and weights depend only on $X_i$ and $X_j$. The sum tested against
each interval is now $W_{ij}$, obtained by deleting both closed
neighborhoods. The dependency graph makes $W_{ij}$ independent of
$(X_i,X_j)$. We can therefore fix the two endpoint values and use the
induction hypothesis on the smaller graph to estimate the interval
probability. Each vertex belongs to at most $q$ closed neighborhoods.
Summing the resulting estimates gives the truncated moment above.
